\IfFileExists{papers/formalized-vs-literature-preamble.tex}{\documentclass[11pt]{article}
\usepackage[utf8]{inputenc}
\usepackage[T1]{fontenc}
\usepackage{amsmath,amssymb,mathtools}
\usepackage{booktabs,tabularx,array}
\usepackage{enumitem}
\usepackage[margin=1in]{geometry}
\usepackage{xcolor}
\usepackage[hidelinks]{hyperref}
\usepackage{microtype}
\setlength{\parindent}{0pt}
\setlength{\parskip}{0.55em}
\newcommand{\Lean}[1]{\texttt{\nolinkurl{#1}}}
\newcommand{\RepoFile}[1]{\texttt{\nolinkurl{#1}}}
\newcommand{\Class}[1]{\textbf{#1}}
\newcommand{\Top}{\mathord{\top}}
\newcommand{\R}{\mathbb{R}}
\newcommand{\ENN}{\mathbb{R}_{\ge 0}^{\infty}}
\newcommand{\KL}{\operatorname{KL}}
\newcommand{\W}{\mathsf{W}}
\newcommand{\statusline}[2]{%
  \begin{center}\fbox{\begin{minipage}{0.92\linewidth}\textbf{Completion class:} #1\\[2pt]
  \textbf{Strongest caveat:} #2\end{minipage}}\end{center}}
\newenvironment{comparison}{\tabularx{\linewidth}{>{\bfseries}p{0.25\linewidth}X}}{\endtabularx}
\newcommand{\snapshot}{\texttt{902f51aa01a737af68f6feb71c83263cc9b1f4d9}}
}{\documentclass[11pt]{article}
\usepackage[utf8]{inputenc}
\usepackage[T1]{fontenc}
\usepackage{amsmath,amssymb,mathtools}
\usepackage{booktabs,tabularx,array}
\usepackage{enumitem}
\usepackage[margin=1in]{geometry}
\usepackage{xcolor}
\usepackage[hidelinks]{hyperref}
\usepackage{microtype}
\setlength{\parindent}{0pt}
\setlength{\parskip}{0.55em}
\newcommand{\Lean}[1]{\texttt{\nolinkurl{#1}}}
\newcommand{\RepoFile}[1]{\texttt{\nolinkurl{#1}}}
\newcommand{\Class}[1]{\textbf{#1}}
\newcommand{\Top}{\mathord{\top}}
\newcommand{\R}{\mathbb{R}}
\newcommand{\ENN}{\mathbb{R}_{\ge 0}^{\infty}}
\newcommand{\KL}{\operatorname{KL}}
\newcommand{\W}{\mathsf{W}}
\newcommand{\statusline}[2]{%
  \begin{center}\fbox{\begin{minipage}{0.92\linewidth}\textbf{Completion class:} #1\\[2pt]
  \textbf{Strongest caveat:} #2\end{minipage}}\end{center}}
\newenvironment{comparison}{\tabularx{\linewidth}{>{\bfseries}p{0.25\linewidth}X}}{\endtabularx}
\newcommand{\snapshot}{\texttt{902f51aa01a737af68f6feb71c83263cc9b1f4d9}}
}
\title{Wang--Gao--Xie Theorem 1: formalized result versus the literature}
\author{}
\date{Formalization snapshot \snapshot}
\begin{document}
\maketitle
\statusline{\Class{R}: genuine Sinkhorn-DRO equality in a strict-interior regime}{The paper's shifted-radius boundary case and $\lambda=0$ endpoint are absent; the current theorem requires a Slater point and positive multipliers.}
\section{Source theorem}
Wang, Gao, and Xie, ``Sinkhorn Distributionally Robust Optimization,'' Theorem~1 treats feasibility, strong duality, finiteness, and optimizer structure using a shifted radius $\bar\rho$.  Its scalar dual includes a limiting $\lambda=0$ term.  The source proof and supplement dependencies are reconstructed in
\RepoFile{prose/distilled_literature/WangGaoXie2025_sinkhorn_dro_theorem1.tex}.

\section{Lean declarations}
The principal equality is \Lean{WangGaoXie2023.strong_duality} in
\RepoFile{WangGaoXie2023/Basic.lean}.  The reverse direction is supplied by
\Lean{WangGaoXie2023.sinkhornDual_le_droValue}, built from the general value-function and Gibbs-attainment machinery in \RepoFile{ForMathlib/OptimalTransport/}.

\section{What Lean proves}
For a standard Borel nonempty space, measurable reward and cost, positive entropy parameter $\kappa$, extensive integrability conditions for the Gibbs kernels, bounded-above reward, a nonempty feasible ball, and a strict feasible budget value $t_0<\epsilon$, Lean proves
\[
 \sup_{\mu\in\mathcal B_{\rm Sink}(\epsilon)}\int f\,d\mu
 =
 \inf_{\lambda>0}\operatorname{DualObjective}(\lambda).
\]
No duality-gap inequality is assumed.  The proof combines edge-free weak duality with concavity of the budget value function and a supergradient multiplier, while the pointwise entropic converse is attained exactly by a Gibbs tilt.

\section{Axis-by-axis comparison}
\begin{comparison}
Radius & Lean exposes the raw objective budget; the source theorem is cleanest in shifted radius $\bar\rho$.\\
Feasibility & Lean proves only the unconditional necessity $\epsilon\ge0$ for a nonempty raw ball.\\
Interior & Lean requires an explicit strict feasible value.\\
Boundary $\bar\rho=0$ & Not proved.\\
Multiplier domain & $\lambda>0$ in Lean; source also includes a $\lambda=0$ limiting endpoint.\\
Value equality & Genuinely proved in the strict-interior regime.\\
Finite value & Encoded indirectly through real-valued definitions and boundedness assumptions.\\
\end{comparison}

\section{Definition of source-level completion}
Introduce the shifted-radius API, prove the exact feasibility equivalence, formalize the essential-supremum endpoint at $\lambda=0$, prove the $\bar\rho=0$ boundary by the source's limiting argument, and then derive the full theorem without Slater.
\end{document}
